\chapter{Conclusion and Future Work} \label{chapter:ch7_conclusion}

This thesis asked whether collaborative-robot tasks that change frequently, represented here by pick-and-place and seam welding, can be authored more directly by specifying spatial intent in the real workspace instead of relying only on lower-level programming workflows such as teach pendant programming, offline programming, or hand-guiding. \DIFdelbegin \DIFdel{Within prototype scope, the answer is positive for selected }\DIFdelend \DIFaddbegin \DIFadd{The implemented system provides a positive answer for prepared }\DIFaddend task-oriented workflows.
\DIFdelbegin \DIFdel{The developed system shows that poses and trajectories entered in the workspace can be combined with scene detection, task-specific interpretation, and robot execution to produce executable procedures on a real robot without exposing the operator to a full robot programming environment}\DIFdelend \DIFaddbegin 

\DIFadd{The main result is that after a one-time setup, calibration, and integration effort, new task instances across both demonstrated task families can be authored through direct spatial input and confirmation inside a use case instead of being rebuilt as low-level robot procedures}\DIFaddend .

\DIFdelbegin \DIFdel{The main contribution is a reusable }\DIFdelend \DIFaddbegin \DIFadd{This is a positive result for prepared }\DIFaddend task-oriented \DIFdelbegin \DIFdel{workflow for spatial authoring, supported by a modular shared runtime with explicit interfaces. Within a selected use case, that workflow supports authoring, confirmation of the interpreted result, execution, visualization, and persistence. The same runtime and interfaces allow both pick-and-place and seam welding to reuse the same infrastructure.
}\DIFdelend \DIFaddbegin \DIFadd{workflows. The rest of the chapter therefore focuses on the next steps that would broaden validation and extend the system toward a wider range of practical conditions.
}\DIFaddend 

\DIFdelbegin \DIFdel{The result nevertheless remains a proof-of-concept. Most development and testing took place in the Kassow REMII offline environment, physical validation is limited, and the thesis does not include a formal user study, a baseline comparison against teach pendant or offline programming workflows, or a quantitative characterization of tracking accuracy, repeatability, or productivity. Within that boundary, seam welding is the clearest practical fit because authored trajectory input maps directly to a welding path.
Pick-and-place mainly shows that an authored pose can be combined with refreshed scene information at execution time, but it is a less compelling practical fit because these tasks are often programmed once and reused rather than re-authored frequently. These limits do not weaken the thesis claim itself; rather, they define the work required to transfer the demonstrated prototype into a broader and more robust system.
}%DIFDELCMD < 

%DIFDELCMD < %%%
\DIFdelend \section{Near-Term Future Work}

Improving sensing, tracking, and calibration is therefore the most direct near-term step for transferring the present result beyond the demonstrated setup. More robust pen tracking, use of pen-side IMU data, and an object-detection pipeline beyond the current fiducial-box representation would directly improve the present workflow. In particular, \DIFdelbegin \DIFdel{the system needs }\DIFdelend \DIFaddbegin \DIFadd{broader applicability would benefit from }\DIFaddend a more accurate and more \DIFdelbegin \DIFdel{broadly applicable }\DIFdelend \DIFaddbegin \DIFadd{generally usable }\DIFaddend way to detect task-relevant objects in the scene, ideally without relying only on predefined tagged objects. \DIFdelbegin \DIFdel{That workshould also }\DIFdelend \DIFaddbegin \DIFadd{These are extension directions beyond the thesis goal rather than missing core deliverables of the present work, and they should }\DIFaddend be accompanied by clearer evaluation of tracking stability and scene-detection accuracy in the real workspace.

\DIFaddbegin \DIFadd{Execution semantics and safety come next. }\DIFaddend Future work should \DIFdelbegin \DIFdel{introduce an explicit abstraction for tool-specific actions, }\DIFdelend \DIFaddbegin \DIFadd{replace the current fixed-delay approximations of gripper and welding actions with an explicit tool-action abstraction. On top of that, broader deployment-oriented extensions such as }\DIFaddend motion planning, collision checking\DIFdelbegin \DIFdel{for authored tasks}\DIFdelend , and simulated execution \DIFdelbegin \DIFdel{for safer inspection before }\DIFdelend \DIFaddbegin \DIFadd{can be added to the current workflow before wider }\DIFaddend real-robot use. \DIFdelbegin \DIFdel{For tasks that do not depend on runtime scene adaptation, it is also worth exploring a path from authored commands to a lower-level robot-side export or compilation target}\DIFdelend \DIFaddbegin \DIFadd{These are natural next steps for broader deployment, not capabilities that the thesis set out to implement}\DIFaddend .

\DIFdelbegin \DIFdel{Separate from those engineering steps, the prototype still needs broader validation and some cross-platform refinements}\DIFdelend \DIFaddbegin \DIFadd{Applicability can be evaluated more systematically through further study}\DIFaddend . Broader physical testing and bounded \DIFdelbegin \DIFdel{characterization of repeatability, tracking accuracy, and }\DIFdelend \DIFaddbegin \DIFadd{quantitative characterization of tracking stability, }\DIFaddend scene-estimation \DIFdelbegin \DIFdel{stability }\DIFdelend \DIFaddbegin \DIFadd{repeatability, and execution consistency }\DIFaddend would make the \DIFdelbegin \DIFdel{present }\DIFdelend \DIFaddbegin \DIFadd{current }\DIFaddend evidence easier to interpret. \DIFdelbegin \DIFdel{Some remaining cross-platform work and compact interface cleanup belong here as secondary tasks, because they affect whether the current prototype can be used reliably in repeated experiments}\DIFdelend \DIFaddbegin \DIFadd{A formal comparison against teach pendant or offline workflows, ideally through a user study focused on frequently changing prepared tasks, would then help determine where this style of spatial authoring is actually advantageous and where conventional workflows remain more practical}\DIFaddend .

\section{Longer-Term Research Directions}

Because the architecture is not tied to the current pen-based interface, richer input modalities remain a natural research direction. Hand tracking, VR or AR controllers, and combinations of speech and gesture could be explored within the same task-oriented model. This would help clarify which forms of spatial input remain natural for different task families while preserving the task-oriented constraints of the current approach.

A longer-term research direction is to study where this style of task-oriented spatial authoring fits best in practice. The thesis establishes the \DIFdelbegin \DIFdel{prototype }\DIFdelend \DIFaddbegin \DIFadd{system }\DIFaddend and points to promising conditions, especially tasks that are re-authored often enough that simpler spatial authoring may matter, but it does not study that practical fit systematically. Future work could examine how spatial authoring should relate to teach pendant or controller workflows and evaluate that question through a formal user study. That would help determine for which task families this approach actually reduces reprogramming effort, where prepared use cases are a good fit, and where more general robot programming remains the better tool.

That research could then guide the next stage of \DIFdelbegin \DIFdel{prototype }\DIFdelend \DIFaddbegin \DIFadd{system }\DIFaddend development and evaluation across a broader range of task families and operating conditions. Richer welding scenarios, additional fabrication or surface-processing tasks, and other human-in-the-loop workflows in which the human provides task intent and the robot performs repeated physical execution would provide a stronger basis for judging that applicability boundary. The present \DIFdelbegin \DIFdel{prototype is a bounded starting point }\DIFdelend \DIFaddbegin \DIFadd{result provides a concrete foundation }\DIFaddend for that work, \DIFdelbegin \DIFdel{but broader claims still depend on technical completion }\DIFdelend \DIFaddbegin \DIFadd{while broader claims about applicability across task families would require corresponding extensions }\DIFaddend and evaluation on \DIFaddbegin \DIFadd{those }\DIFaddend concrete task families.
